\documentclass[11pt,a4paper]{amsart}

\usepackage[T1]{fontenc}
\usepackage[utf8]{inputenc}
\usepackage[a4paper,margin=1in]{geometry}
\usepackage{amsmath,amssymb,mathtools}
\usepackage{microtype}
\usepackage[hidelinks]{hyperref}
\emergencystretch=2em

\newtheorem{theorem}{Theorem}[section]
\newtheorem{proposition}[theorem]{Proposition}
\newtheorem{lemma}[theorem]{Lemma}
\newtheorem{corollary}[theorem]{Corollary}
\theoremstyle{definition}
\newtheorem{definition}[theorem]{Definition}
\theoremstyle{remark}
\newtheorem{remark}[theorem]{Remark}

\newcommand{\R}{\mathbb R}
\newcommand{\Abar}{\overline A}
\newcommand{\cross}{\mathbin{\times}}

\title[An exact zero of focal antipedal areas]{An Exact Zero of Focal Antipedal Areas}
\author{DannyExperiments}
\date{August 12, 2026}
\subjclass[2020]{37J35, 51M04}
\keywords{elliptic billiard, Poncelet polygon, antipedal polygon, signed area, exact counterexample}
\hypersetup{
  pdftitle={An Exact Zero of Focal Antipedal Areas},
  pdfauthor={DannyExperiments},
  pdfsubject={An exact domain counterexample to arXiv:2004.12497v11, row k607}
}

\begin{document}

\begin{abstract}
Row $k_{607}$ of arXiv:2004.12497v11 proposes that the quotient of the two focal antipedal signed areas is $1$ for elliptic-billiard periods divisible by four.
Under the standard complete-supporting-line construction, we prove that central inversion equates the two signed areas whenever the orbit has the corresponding half-period symmetry.
We then give an exact primitive one-winding $8$-periodic orbit for which all sixteen focal antipedal intersections are finite but both signed areas vanish.
The displayed quotient is therefore $0/0$ at that admissible phase, so the unqualified quotient statement needs a nonzero-denominator hypothesis, while the signed-area identity survives.
The finite algebraic certificate is independently reproducible, but the geometric and domain arguments are proved analytically here.
The exact $N=8$ certificate appears new after a documented three-lane search through August 12, 2026; absolute historical priority is not claimed.
\end{abstract}

\maketitle

\section{The problem and the answer}

Let
\[
 E:\frac{x^2}{a^2}+\frac{y^2}{b^2}=1,
 \qquad a>b>0,
\]
have foci $F_+=(c,0)$ and $F_-=(-c,0)$, where
$c^2=a^2-b^2$.  Let $P=(P_0,\ldots,P_{N-1})$ be an ordered periodic
elliptic-billiard orbit tangent to a fixed nondegenerate confocal caustic.
For a point $F$ and a vertex $P_i$, let $L_i(F)$ denote the complete line
through $P_i$ perpendicular to $P_i-F$.  Whenever consecutive lines are
nonparallel, put
\[
 Q_i(F)=L_i(F)\cap L_{i+1}(F),
 \qquad
 \Abar^*(F)=\frac12\sum_{i\bmod N}Q_i(F)\cross Q_{i+1}(F),
\]
where $(x_1,y_1)\cross(x_2,y_2)=x_1y_2-y_1x_2$.  Thus the area is the
cyclic signed shoelace area; self-intersections do not trigger absolute
values.

Reznik, Garcia, and Koiller define polygonal areas with this signed
convention and list, in row $k_{607}$ of the source table of
arXiv:2004.12497v11, the proposed value
\begin{equation}\label{eq:source-problem}
 \frac{\Abar^*(F_+)}{\Abar^*(F_-)}=1,
 \qquad N\equiv0\pmod 4,
\end{equation}
with proof status ``?'' \cite{ReznikGarciaKoiller2020v11}.  That exact row was
first present in this form in the revision of October 29, 2020.  This paper is
about \emph{that versioned formula}.  The later journal article uses the
identifier $k607$ for a different pedal identity and omits
\eqref{eq:source-problem}; it treats analogous focal-antipedal relations as
expected but unchecked \cite[Sec.~3.7 and Table~7]{ReznikGarciaKoiller2021}.

Our answer has two parts.

\begin{theorem}[Corrected focus-exchange statement]\label{thm:corrected}
Let $N=2m$ and suppose an ordered polygon satisfies
$P_{i+m}=-P_i$.  For the opposite points $F_-=-F_+$, assume that all
consecutive complete supporting-line intersections $Q_i(F_+)$ and
$Q_i(F_-)$ are finite.  Then
\[
 Q_{i+m}(F_-)=-Q_i(F_+)
 \quad\text{and}\quad
 \Abar^*(F_-)=\Abar^*(F_+).
\]
Consequently the quotient in \eqref{eq:source-problem} equals $1$ wherever
the common signed area is nonzero and is undefined wherever that area is
zero.
\end{theorem}

\begin{theorem}[Exact $N=8$ domain counterexample]\label{thm:counterexample}
There exist a noncircular ellipse, a strictly interior nondegenerate confocal
elliptic caustic, and a primitive counterclockwise one-winding $8$-periodic
billiard orbit such that every side is tangent to the caustic on its open
segment, all sixteen focal antipedal supporting-line intersections are finite,
and
\[
 \Abar^*(F_+)=\Abar^*(F_-)=0.
\]
Thus the ordinary real quotient in \eqref{eq:source-problem} is undefined at
this admissible phase.
\end{theorem}

Theorem~\ref{thm:counterexample} is a denominator or domain counterexample,
not an example with unequal focal areas.  The underlying equality in
Theorem~\ref{thm:corrected} remains valid.  The exact $N=8$ certificate
appears new after a documented three-lane search through August 12, 2026,
with moderate-high bounded-search confidence; that search is not a proof of
absolute priority.

The supporting-line convention is deliberate.  Classical antipedal geometry
uses complete perpendicular lines \cite[Chap.~VII]{Gallatly1910}; modern
treatments retain the same convention \cite{Mitrea2010,GarciaEtAl2020Hat}.
The source prose says ``rays'' without specifying directions, while its figure
and public implementation use unrestricted affine-line intersections
\cite{ReznikGarciaKoiller2020v11,ReznikEtAlRepository2020}.  We make no theorem about a polygon formed
from one selected half-ray at each vertex.  For the witness below such a
half-ray cycle does not form; see Remark~\ref{rem:rays}.

The ingredients surrounding the result are not new.  Central symmetry of
simple even elliptic billiards of odd turning number is known
\cite[Cor.~4.2(i)]{Stachel2021}, and the source authors already use opposite
vertices to prove the analogous focal-pedal equality
\cite[Sec.~3.7]{ReznikGarciaKoiller2021}.  Zero antipedal areas were also
known computationally in other settings \cite{Reznik2020Regular}.  More
specifically, on October 28, 2020 Reznik publicly exhibited an experimental
$N=6$, $a/b=2$ elliptic-billiard family whose left-focus antipedal signed area
was dynamically zero \cite{Reznik2020N6Focus}.  That example lies outside
$N\equiv0\pmod4$ and does not contain the exact simultaneous-two-focus $N=8$
certificate below.  Accordingly, no novelty claim is made for complete-line
antipedals, central symmetry, the focus-exchange mechanism, or zero antipedal
phenomena in general.

Exact algebraic coordinates for a generic simple symmetric axis-position
$N=8$ orbit also predate the present certificate
\cite[Appendix~B.7.1]{GarciaReznik2022}.  Their parameter satisfies a
different quartic involving the ellipse data.  We therefore claim neither
the first algebraization of such octagons nor a new generic coordinate
architecture; the bounded novelty claim concerns only the particular quartic
specialization below and its simultaneous focal zero.

\section{Central inversion and the corrected theorem}

Let $R(x,y)=(-y,x)$ denote counterclockwise rotation through $\pi/2$.
The complete antipedal supporting line is parametrized by
\begin{equation}\label{eq:line-param}
 L_i(F)=\{P_i+tR(P_i-F):t\in\R\}.
\end{equation}

\begin{proof}[Proof of Theorem~\ref{thm:corrected}]
Let $J(X)=-X$.  From $P_{i+m}=-P_i$, $F_-=-F_+$, and linearity of $R$,
\[
 J\bigl(P_i+tR(P_i-F_+)\bigr)
 =P_{i+m}+tR(P_{i+m}-F_-).
\]
Hence $J(L_i(F_+))=L_{i+m}(F_-)$.  Since consecutive intersections are
unique,
\[
 Q_{i+m}(F_-)=-Q_i(F_+).
\]
The index change $i\mapsto i+m$ is a cyclic shift, not a reversal.  Moreover
$\det(-I_2)=1$, and therefore $(-X)\cross(-Y)=X\cross Y$.  Reindexing the
shoelace sum gives
\[
 \begin{aligned}
 2\Abar^*(F_-)
 &=\sum_{i\bmod N}Q_{i+m}(F_-)\cross Q_{i+m+1}(F_-)\\
 &=\sum_{i\bmod N}(-Q_i(F_+))\cross(-Q_{i+1}(F_+))
 =2\Abar^*(F_+).
 \end{aligned}
\]
The statement about the quotient is now ordinary real division with its
nonzero-denominator condition retained.
\end{proof}

For the simple one-winding even elliptic-caustic billiards relevant below,
the half-period relation $P_{i+N/2}=-P_i$ is standard
\cite[Cor.~4.2(i)]{Stachel2021}.  Thus Theorem~\ref{thm:corrected} applies to
that established scope.  We use the cited central-symmetry theorem as prior
input and make no claim of a new central-symmetry result or of an extension to
other caustic or turning classes.

We also record that the finiteness hypothesis is automatic for a
nondegenerate confocal billiard when complete supporting lines are used.

\begin{lemma}\label{lem:finite}
No real tangent to a nondegenerate confocal caustic passes through a focus of
the caustic.  Consequently consecutive focal antipedal supporting lines of a
nondegenerate confocal billiard orbit are nonparallel.
\end{lemma}

\begin{proof}
Write the caustic as
\[
 C_\lambda:\frac{x^2}{A}+\frac{y^2}{B}=1,
 \qquad A>0,\quad B\ne0,
\]
so its foci are $(\pm c,0)$ with $c^2=A-B$.  A line
$ux+vy=w$ is tangent precisely when
\[
 w^2=Au^2+Bv^2.
\]
If it passed through $(c,0)$, then $w=uc$ and substitution would give
$B(u^2+v^2)=0$, impossible for a line.  The argument for $(-c,0)$ is
identical.

Now $L_i(F)$ and $L_{i+1}(F)$ are parallel exactly when
$(P_i-F)\cross(P_{i+1}-F)=0$, equivalently when the billiard side
$P_iP_{i+1}$ passes through $F$.  Since that side is tangent to the caustic,
the first part excludes this possibility.
\end{proof}

\section{The exact octagon}

We now prove Theorem~\ref{thm:counterexample}.  Let $u$ be the real root in
\begin{equation}\label{eq:root-interval}
 \frac{313}{1000}<u<\frac{157}{500}
\end{equation}
of
\begin{equation}\label{eq:g}
 g(u)=u^4-6u^3-2u^2-2u+1.
\end{equation}
The exact endpoint values are
\[
 g(313/1000)=\frac{3674142961}{10^{12}}>0,
 \qquad
 g(157/500)=-\frac{76605799}{62500000000}<0.
\]
Also
\[
 g'(u)=4u(u^2-1)-18u^2-2<0\qquad(0<u<1),
\]
so the root exists and is unique in the stated interval.  In particular
$u<157/500<1/3<\sqrt2-1$.

Set
\begin{equation}\label{eq:parameters}
 r=\frac{(1-u)^3(1+u)}{4u^3},\qquad
 a=\sqrt r,\qquad b=1,\qquad c=\sqrt{r-1},
\end{equation}
and
\begin{equation}\label{eq:CS}
 C=\frac{1-u^2}{1+u^2},\qquad
 S=\frac{2u}{1+u^2}.
\end{equation}
The factorization
\begin{equation}\label{eq:rminusone}
 r-1=-\frac{(1+u^2)(u^2+2u-1)}{4u^3}>0
\end{equation}
shows $a>1$.  Define
\begin{equation}\label{eq:vertices}
\begin{aligned}
P_0&=(a,0),&P_1&=(aC,S),&P_2&=(0,1),&P_3&=(-aC,S),\\
P_4&=(-a,0),&P_5&=(-aC,-S),&P_6&=(0,-1),&P_7&=(aC,-S).
\end{aligned}
\end{equation}
Because $C^2+S^2=1$, all vertices lie on
\begin{equation}\label{eq:outer}
 E:\frac{x^2}{a^2}+y^2=1.
\end{equation}
Let $\theta=2\arctan u$.  From $0<u<\sqrt2-1$ we have
$0<\theta<\pi/4$, and the eccentric anomalies of the displayed vertices are
\[
 0,\ \theta,\ \frac\pi2,\ \pi-\theta,\ \pi,\quad
 \pi+\theta,\ \frac{3\pi}2,\ 2\pi-\theta.
\]
They are strictly increasing modulo $2\pi$.  Thus the vertices are distinct,
occur counterclockwise, complete one winding, and satisfy $P_{i+4}=-P_i$.
In particular, no lower-period traversal is repeated.

\subsection*{One common caustic and the billiard law}

Put
\begin{equation}\label{eq:lambda}
 \lambda=\frac{a^2u^2}{1+a^2u^2}.
\end{equation}
Then $0<\lambda<1$, so
\begin{equation}\label{eq:caustic}
 C_\lambda:\frac{x^2}{a^2-\lambda}
             +\frac{y^2}{1-\lambda}=1
\end{equation}
is a nondegenerate confocal ellipse strictly inside $E$.

For $P(\alpha)=(a\cos\alpha,\sin\alpha)$, the line through
$P(\alpha)$ and $P(\beta)$ has equation
\[
 \frac{x}{a}\cos\mu+y\sin\mu=\cos\delta,
 \qquad
 \mu=\frac{\alpha+\beta}{2},\quad
 \delta=\frac{\alpha-\beta}{2}.
\]
The dual-conic tangency criterion gives its confocal parameter as
\begin{equation}\label{eq:chord-lambda}
 \lambda(\alpha,\beta)
 =\frac{\sin^2\delta}{\cos^2\mu/a^2+\sin^2\mu}.
\end{equation}
For the two representative side types, formula \eqref{eq:chord-lambda}
becomes
\[
 \lambda_{01}=\frac{a^2u^2}{1+a^2u^2},
 \qquad
 \lambda_{12}=\frac{(1-u)^2}{(1-u)^2/a^2+(1+u)^2}.
\]
Their equality is equivalent to
\begin{equation}\label{eq:closure}
 h(u,a)=u^4+(4a^2-2)u^3+2u-1=0.
\end{equation}
But \eqref{eq:parameters} says
$4a^2u^3=(1-u)^3(1+u)$, and expanding this identity yields
\eqref{eq:closure}.  Reflections in the coordinate axes carry the two
representative sides to all eight sides, so every side is tangent to the one
caustic \eqref{eq:caustic}.  Since the caustic is strictly inside the convex
ellipse $E$, the contact on each chord line lies between its two intersections
with $E$, hence on the open side segment.

For completeness, the reflection law can be checked without invoking the
general confocal correspondence.  At $P_1=(aC,S)$, equality of the tangential
components of the incoming and outgoing unit directions is
\begin{equation}\label{eq:reflection-unsquared}
 \frac{2a^2u^2+1-u^2}{\sqrt{a^2u^2+1}}
 =\frac{(1+u)\bigl(2a^2u+(1-u)^2\bigr)}
 {\sqrt{a^2(1+u)^2+(1-u)^2}}.
\end{equation}
Both sides are positive.  After squaring and clearing the positive
denominators, left square minus right square is
\[
 -a^2(1+u^2)^2h(u,a)=0.
\]
Thus \eqref{eq:reflection-unsquared} holds without an extraneous sign.  The
axis vertices satisfy reflection by mirror symmetry, and coordinate
reflections cover every remaining vertex.  The octagon is therefore an exact
primitive one-winding billiard orbit with a single nondegenerate caustic.

\subsection*{Finite antipedal vertices and zero signed area}

Take $F_+=(c,0)$ and set
\[
 D_i=(P_i-F_+)\cross(P_{i+1}-F_+).
\]
The four representative determinants are
\begin{align}
D_0&=\frac{2u(a-c)}{1+u^2},
&D_1&=\frac{(1-u)\bigl(a(1+u)-c(1-u)\bigr)}{1+u^2},\label{eq:dets}\\
D_2&=\frac{(1-u)\bigl(a(1+u)+c(1-u)\bigr)}{1+u^2},
&D_3&=\frac{2u(a+c)}{1+u^2}.
\end{align}
The remaining four occur in reverse reflected order.  Because
$0<u<1$ and $a>c>0$, every determinant in \eqref{eq:dets} is strictly
positive.  Hence all eight plus-focus intersections are unique and finite;
central inversion gives the same conclusion for the minus focus.

We give enough algebra to make the area certificate directly checkable.  For
$n_i=P_i-F_+$ and $s_i=P_i\cdot n_i$, the line at $P_i$ has equation
$X\cdot n_i=s_i$.  Solving two consecutive equations gives
\begin{equation}\label{eq:intersection}
 Q_i(F_+)=
 \frac{s_{i+1}R(n_i)-s_iR(n_{i+1})}{n_i\cross n_{i+1}}.
\end{equation}
Substitute \eqref{eq:vertices} into \eqref{eq:intersection}, form the cyclic
shoelace sum, and reduce using only
\begin{equation}\label{eq:relations}
 c^2=a^2-1,
 \qquad 4a^2u^3=(1-u)^3(1+u).
\end{equation}
The exact result is
\begin{equation}\label{eq:area-factor}
 \Abar^*(F_+)
 =\frac{a(u^2-2u-1)
 \bigl(u^4-6u^3-2u^2-2u+1\bigr)}{4u^3}.
\end{equation}
For an explicit check that no forbidden cancellation creates the zero, before
the reduction the rational denominator contains the strictly positive factor
\[
 (1+u^2)^2\bigl(4a^2u+(1-u)^2\bigr).
\]
Equation \eqref{eq:g} now makes \eqref{eq:area-factor} zero.  Finally,
$P_{i+4}=-P_i$ and Theorem~\ref{thm:corrected} give
\[
 \Abar^*(F_-)=\Abar^*(F_+)=0.
\]
All intersections are finite, so the zero comes from signed shoelace
cancellation in a self-overlapping polygon, not from a point at infinity or
an undefined intermediate construction.  The quotient in
\eqref{eq:source-problem} is exactly $0/0$.  This proves
Theorem~\ref{thm:counterexample}.

\begin{remark}[Why this does not prove a half-ray theorem]\label{rem:rays}
The argument uses the complete lines \eqref{eq:line-param}.  If one instead
selects, at every vertex, the half-ray with parameter $t\ge0$ (or reverses
all of them), then at lines $0,1,4,7$ the two adjacent supporting-line
intersections have opposite parameter signs.  No selected half-ray there
contains both required adjacent vertices, so the cyclic antipedal polygon is
not formed.  Central inversion preserves $t$ when a ray point exists; it does
not create missing intersections.  Thus the source's unoriented word
``rays'' requires a separate definitional repair and is outside the theorem.
\end{remark}

\section{Verification, priority, and limits}

The proof above is analytic and exact.  Root isolation uses rational endpoint
signs and monotonicity; cyclic order uses the isolating interval; caustic
tangency and reflection use \eqref{eq:closure}; finiteness uses the positive
determinants \eqref{eq:dets}; and the decisive zero is the factorization
\eqref{eq:area-factor}.  Those arguments, rather than software output, carry
the geometric theorem.

As corroboration, a dependency-free standard-library program performs exact
arithmetic in the quartic extension and verifies all eight outer-ellipse
relations, all eight dual-caustic tangency equations, all sixteen finite
supporting-line intersections, the central-inversion vertex map, and both
zero signed areas.  It replays identically under ordinary and optimized
Python.  A separate partial Lean~4 artifact kernel-checks the finite algebraic
certificate but not the real-root embedding, positivity, open-segment
geometry, unsquared reflection sign, cyclic order, primitivity, source
semantics, or priority.  Accordingly, this paper does not claim a formal proof
of the complete theorem.

The literature boundary is equally specific.  Complete-line antipedal
geometry and triangle area relations are classical
\cite{Gallatly1910,Mitrea2010}.  Central symmetry supplies the focus-exchange
mechanism \cite{Stachel2021}, and the source authors' focal-pedal proof is its
closest published analogue \cite{ReznikGarciaKoiller2021}.  Their public media
also show both regular-polygon zero-area phenomena and the experimental
one-focus $N=6$, $a/b=2$ focal example just described
\cite{Reznik2020Regular,Reznik2020N6Focus,ReznikMediaIndex}.
No exact match was retrieved, through August 12, 2026, for the primitive
noncircular focal $N=8$ simultaneous-zero certificate, its quartic, and the
complete exact verification.  The accurate novelty statement is therefore
only that this exact certificate \emph{appears new after that documented
search}.

No claim is made about a literal half-ray construction, the differently
numbered journal item, unsigned area, odd periods, other-period zero
existence, other caustics, a classification or uniqueness of the common-zero
locus, or a second open problem.  Nor is the quotient assigned a removable
value at $0/0$; doing so would define a regularized quantity different from
the pointwise real quotient in \eqref{eq:source-problem}.

\section*{Acknowledgments and computational disclosure}

DannyExperiments initiated and directed the investigation, selected the
problem, curated the evidence, and is the sole manuscript author.  An OpenAI
ChatGPT Pro instance produced the frozen proof candidate; OpenAI Codex assisted with
independent reconstruction, mathematical and literature audits, exact
verification, limited formalization, and manuscript preparation.  No AI
system is listed as an author.  These AI-assisted audits are not human peer
review; responsibility for the manuscript remains with the named author.

\bibliographystyle{amsplain}
\bibliography{references}

\end{document}
